% GENERATED by scripts/build-lecture-slides.py; do not edit. Portable profile: XeTeX, theme and fonts from the core.
\makeatletter\IfFileExists{core/beamer/beamerthememx.sty}{\def\mxcorepath{core/}\def\input@path{{core/beamer/}{core/}}}{\def\mxcorepath{../core/}\def\input@path{{../core/beamer/}{../core/}}}\makeatother
\documentclass[aspectratio=169,10pt,t]{beamer}
\usetheme{mx}
\usepackage{graphicx,booktabs}
\usepackage[french]{babel}
\setlength{\parskip}{.45em}
\hypersetup{pdfauthor={},pdfsubject={Lecture slides},pdfkeywords={},pdftitle={Systèmes d’équations linéaires dans R\textsuperscript{3}}}
\providecommand{\N}{\mathbb{N}}
\providecommand{\Z}{\mathbb{Z}}
\providecommand{\D}{\mathbb{D}}
\providecommand{\Q}{\mathbb{Q}}
\providecommand{\R}{\mathbb{R}}
\providecommand{\C}{\mathbb{C}}
\providecommand{\vect}[1]{\overrightarrow{#1}}
\providecommand{\vi}{\vec{\imath}}
\providecommand{\vj}{\vec{\jmath}}
\providecommand{\vk}{\vec{k}}
\providecommand{\norm}[1]{\left\lVert #1\right\rVert}
\providecommand{\vc}[2]{\begin{pmatrix}#1\\#2\end{pmatrix}}
\providecommand{\deter}[2]{\det\left(#1,#2\right)}
\providecommand{\abs}[1]{\left\lvert #1\right\rvert}
\providecommand{\dd}{\,\mathrm{d}}
\providecommand{\e}{\mathrm{e}}
\providecommand{\Card}{\operatorname{Card}}
\providecommand{\E}{\operatorname{E}}
\providecommand{\V}{\operatorname{V}}
\providecommand{\degre}{^{\circ}}
\providecommand{\un}[1]{\,\mathrm{#1}}
\providecommand{\sci}[2]{#1\times10^{#2}}
\providecommand{\pgcd}{\operatorname{PGCD}}
\providecommand{\ppcm}{\operatorname{PPCM}}
\title{Systèmes d’équations linéaires dans R\textsuperscript{3}}
\subtitle{Cours et exemples}
\date{}
\author{}
\begin{document}
\mxtitleframe
\begin{frame}{Introduction}
On s'intéresse à la résolution de systèmes de trois équations linéaires à trois inconnues ;
c'est-à-dire à des systèmes de la forme :
\[
	\left\{\begin{array}{l}
		a_1 x + b_1 y + c_1 z = d_1	\\
		a_2 x + b_2 y + c_2 z = d_2	\\
		a_3 x + b_3 y + c_3 z = d_3	.
	\end{array}\right.
\]

Ce genre de système peut aisément se généraliser en système de \(n\) équations à \(n\) inconnues, 
la méthode de résolution étant essentiellement similaire.
\end{frame}
\begin{frame}{Résolution de systèmes linéaires}
Considérons le système suivant :
                              \[
                                    \left\lbrace
                                    \begin{array}{lc}
                                          3 x + 2 y -   z =   1  & (1) \\
                                          2 x - 3 y +   z =  10  & (2) \\
                                          - x +   y - 3 z = -12. & (3)
                                    \end{array}
                                    \right.
                              \]
\end{frame}
\begin{frame}{Résolution de systèmes linéaires (2)}
On choisit une des équations pour être le \emph{pivot}.
\end{frame}
\begin{frame}{Résolution de systèmes linéaires (3)}
On va prendre l'équation \((1)\).
                              \[
                                    \left\lbrace
	                                    \begin{array}{lcl}
	                                          3 x + 2 y -   z =   1  
	                                                & (1) 
	                                                      & \leftarrow \text{pivot} \\
	                                          2 x - 3 y +   z =  10  & (2) &            \\
	                                          - x +   y - 3 z = -12. & (3) &
	                                    \end{array}
                                    \right.
                              \]
\end{frame}
\begin{frame}{Résolution de systèmes linéaires}
On utilise le pivot pour éliminer une des variables dans les autres équations :
                              \[
                                    \begin{matrix}
		                                          +
		                                    \underline{
			                                    \left\lbrace
				                                    \begin{array}{lc}
				                                          3 x + 2 y -   z =   1  & (1) \\
				                                          2 x - 3 y +   z =  10  & (2) 
				                                    \end{array}
			                                    \right.
			                              }
			                              \\
	                                                5 x - y = 11 \quad (2')
	                              \end{matrix}
                                    \qquad
                                    \begin{matrix}
		                                          -
		                                    \underline{
			                                    \left\lbrace
				                                    \begin{array}{lc}
				                                          9 x + 6 y - 3 z =   3  & 3\times(1) \\
				                                          - x +   y - 3 z = -12. & (3)
				                                    \end{array}
			                                    \right.                                    
			                              }
			                              \\
			                              10 x + 5 y = 15 \quad (3')
                                    \end{matrix}
                              \]
\end{frame}
\begin{frame}{Résolution de systèmes linéaires (2)}
Ce qui nous donne le nouveau système :
                              \[
                                    \left\lbrace
	                                    \begin{array}{lcl}
	                                           3 x + 2 y -   z =   1  
	                                                & (1) 
	                                                      & \leftarrow \text{on ne touche plus à l'ancien pivot,}  \\
	                                          5 x - y         =  11  & (2') & \\
	                                          10 x + 5 y      =  15. & (3') & \\
	                                    \end{array}
                                    \right.
                              \]
\end{frame}
\begin{frame}{Résolution de systèmes linéaires}
On choisit alors un nouveau pivot, prenons l'équation \((2')\).
                              \[
                                    \left\lbrace
	                                    \begin{array}{lcl}
	                                          3 x + 2 y -   z =   1  & (1) &            \\
	                                          5 x - y = 11  
	                                                & (2') 
	                                                      & \leftarrow \text{pivot} \\
	                                          10 x + 5 y = 15 & (3') &
	                                    \end{array}
                                    \right.
                              \]
\end{frame}
\begin{frame}{Résolution de systèmes linéaires (2)}
A l'aide de ce nouveau pivot, on élimine une variable dans toutes les équations,
                              excepté celles qui ont déjà servi de pivots.
                              \[
                                    \begin{matrix}
	                                          +
                                          \underline{
		                                    \left\lbrace
			                                    \begin{array}{lcl}
			                                          25 x - 5 y = 55 & 5\times(2') & \\
			                                          10 x + 5 y = 15 & (3') &  
			                                    \end{array}
		                                    \right.
		                                    }
                                                \\
                                                            35 x        = 70      \quad (3'')
                                    \end{matrix}
                              \]
\end{frame}
\begin{frame}{Résolution de systèmes linéaires (3)}
On a alors le système suivant : 
                              \[
                                    \left\lbrace
	                                    \begin{array}{lc}
	                                           3 x + 2 y -   z =   1  & (1 )  \\
	                                           5 x -   y       =  11  & (2') \\
	                                          35 x             =  70. & (3") \\
	                                    \end{array}
                                    \right.
                              \]
\end{frame}
\begin{frame}{Résolution de systèmes linéaires}
On résout les équations une par une :
                              \[
                                    \left\lbrace
	                                    \begin{array}{lc}
	                                          35 x             =  70. & (3") \\
	                                           5 x -   y       =  11  & (2') \\
	                                           3 x + 2 y -   z =   1  & (1 )  \\
	                                    \end{array}
                                    \right.
                                    \Rightarrow
                                    \left\lbrace
	                                    \begin{array}{l}
	                                             x                   =   2    \\
	                                           5 \cdot 2 -   y       =  11    \\
	                                           3 \cdot 2 + 2 y -   z =   1    \\
	                                    \end{array}
                                    \right.
                                    \Rightarrow
                                    \left\lbrace
	                                    \begin{array}{l}
	                                             x           =  2  \\
	                                           10 -  y       =  11 \\
	                                           6 + 2 y -   z =   1 
	                                    \end{array}
                                    \right.
                              \]
\end{frame}
\begin{frame}{Résolution de systèmes linéaires (2)}
\[
                                    \Rightarrow
                                    \left\lbrace
	                                    \begin{array}{l}
	                                           x =  2       \\
	                                           y = -1       \\
	                                           2 \cdot (-1) -   z =   1 -6 
	                                    \end{array}
                                    \right.
                                    \Rightarrow
                                    \left\lbrace
	                                    \begin{array}{l}
	                                           x =  2 \\
	                                           y = -1 \\
	                                           - 2 - z = -5
	                                    \end{array}
                                    \right.
                                    \Rightarrow
                                    \left\lbrace
	                                    \begin{array}{l}
	                                           x =  2 \\
	                                           y = -1 \\
	                                           z =  3 .
	                                    \end{array}
                                    \right.
                              \]
\end{frame}
\begin{frame}{Résolution de systèmes linéaires}
Ainsi, \((x,y,z) = (2,-1,3)\).
\end{frame}
\begin{frame}{Résolution de systèmes linéaires}
Considérons le système suivant :
				\[
					\left\lbrace
						\begin{array}{lc}
							  2 x - 3 y + z = -7  & (1) \\
							  3 x +   y - z = 11  & (2) \\
							- 2 x + 2 y + z =  0. & (3)
						\end{array}
					\right.
				\]
\end{frame}
\begin{frame}{Résolution de systèmes linéaires (2)}
On choisit une équation et une variable à substituer.
			Dans notre cas, l'équation \((2)\) donne \(z = 3 x + y - 11\).
			On a alors 
				\[
					\left\lbrace
						\begin{array}{l}
							  2 x - 3 y + 3 x + y - 11 = -7  \\
							- 2 x + 2 y + 3 x + y - 11 =  0. 
						\end{array}
					\right.
					\Rightarrow
					\left\lbrace
						\begin{array}{lc}
							5 x - 2 y =  4  & (1') \\
							  x + 3 y = 11. & (3')
						\end{array}
					\right.
				\]

(Ce système peut être résolu par n'importe quelle méthode.)
\end{frame}
\begin{frame}{Résolution de systèmes linéaires}
On utilise l'équation \((3')\) pour avoir \(x = -3y+11\).
			Ceci donne 
			      \[ 
			            5 \cdot (-3y+11) - 2 y = 4
			            \Rightarrow
			            -15 y + 55 - 2y = 4
			            \Rightarrow
			            55 - 4 = 17 y .
			      \]

On a alors :

\begin{itemize}

			      \item \(y = \dfrac{51}{17} = 3\),
			      
\item \(x = -3 \cdot 3 + 11 = 2\),
			      
\item \(z = 3 \cdot 2 + 3 - 11 = -2\).

\end{itemize}

			La solution est donc \((x,y,z) = (2,3,-2)\).
\end{frame}
\begin{frame}{Résolution de systèmes linéaires}
Un système linéaire de trois équations à trois inconnues a soit une solution unique,
            soit aucune solution (on dit qu'il est \emph{impossible}), soit une infinité de solutions (on dit qu'il est \emph{indéterminé}).
            La méthode de Gauss permet de le savoir : si, au cours des éliminations, on obtient une égalité
\end{frame}
\begin{frame}{Résolution de systèmes linéaires (2)}
\begin{itemize}

                  \item fausse, du type \(0 = 5\), le système n'a aucune solution ;
                  
\item toujours vraie, du type \(0 = 0\), une équation est conséquence des autres : on garde les équations restantes,
                        on choisit une inconnue comme paramètre et on exprime les autres en fonction de ce paramètre.

\end{itemize}
\end{frame}
\begin{frame}{Résolution de systèmes linéaires}
Considérons
                  \(\left\lbrace
                        \begin{array}{lc}
                              x + y + z = 2 & (1) \\
                              x - y + 2 z = 1 & (2) \\
                              2 x + 3 z = 3 & (3)
                        \end{array}
                  \right.\).
                  En additionnant \((1)\) et \((2)\), on obtient \(2x + 3z = 3\), qui est exactement \((3)\) :
                  l'équation \((3)\) n'apporte rien. On pose \(z = t\) (paramètre réel) ; \((3)\) donne \(x = \dfrac{3 - 3t}{2}\),
                  puis \((1)\) donne \(y = 2 - x - t = \dfrac{1 + t}{2}\).
                  L'ensemble des solutions est
                  \(\left\lbrace \left( \dfrac{3 - 3t}{2} ; \dfrac{1 + t}{2} ; t \right) \mid t \in \mathbb{R} \right\rbrace\) :
                  il y en a une infinité.
                  
                  Si la troisième équation était \(2x + 3z = 4\), on obtiendrait \(3 = 4\) : le système n'aurait aucune solution.
\end{frame}
\end{document}
